\chapter{Glossary}
\label{ch:glossary}

\begin{longtable}{L{0.26\textwidth}L{0.66\textwidth}}
	\caption{Glossary}\label{tab:glossary}\\
	\toprule
	\textbf{Term} & \textbf{Definition} \\
	\midrule
	\endfirsthead
	\toprule
	\textbf{Term} & \textbf{Definition} \\
	\midrule
	\endhead

	Abstraction & The deliberate discarding of code constructs that sit below the
	level at which a systems engineer reads a model. \\
	\addlinespace
	Activity & A sequential behaviour model of actions, decisions, forks and
	merges. \\
	\addlinespace
	Control flow graph & Basic blocks and the edges between them, produced from
	source by the compiler platform. \\
	\addlinespace
	Entity & A record in the GENESYS repository. Diagrams in GENESYS are views
	generated from entities and relationships rather than stored pictures. \\
	\addlinespace
	Intermediate representation & The language-neutral model exchanged between the
	analysis and mapping halves of the tool. \\
	\addlinespace
	Semantic model & The compiler's resolved view of a solution, which maps a
	syntax node to the symbol it actually references. \\
	\addlinespace
	State machine & A set of states, transitions, guards, events and entry or exit
	behaviours. \\
	\addlinespace
	Translation & Either direction of conversion between a GENESYS model and
	source code. \\
	\bottomrule
\end{longtable}
